\section{Warm-up: linear-regression experts}

With linear experts, EM has closed-form steps, so every part of the machinery (complete-data likelihood, E-step, M-step) can be derived by hand. Sections 3--5 then swap the linear experts for MLPs and the softmax gate for an HMM. The structure of the derivation stays the same; only the M-step loses its closed form.

\subsection{Reading guide (about 2.5 hours)}

The main reference is F.~Chamroukhi and H.~D.~Nguyen, \emph{Model-Based Clustering and Classification of Functional Data}, arXiv:1803.00276 (2018). Page numbers below are the printed ones.

The paper has \emph{two} layers of latent variables:
\begin{itemize}
  \item an outer cluster label $Z_i$ for a whole curve $i$;
  \item an inner regime $H_{ij}$ at each point $j$ of the curve.
\end{itemize}
Keep this in mind while reading.

\begin{enumerate}[label=\textbf{R\arabic*.}]
\item \textbf{\S2.1--2.2, pp.\ 8--11 (20 min).} The general mixture, the complete-data log-likelihood (eq.~4), and EM in general form (eqs.~5--8). This is the template that every later section fills in.
\item \textbf{\S3.1--3.2, pp.\ 14--16 (30 min).} A mixture of linear regressions and its EM. The M-step is weighted least squares (eqs.~18--19).
  \begin{itemize}
    \item Each curve $i$ has \emph{one} label $Z_i$ shared by all of its $m_i$ points, so the component density is a product over those points.
    \item This is exactly our ``one regime per date, shared by every stock'': read \emph{curve} as \emph{date} and \emph{point} as \emph{stock}.
  \end{itemize}
\item \textbf{\S4.3.1--4.3.2, pp.\ 48--51 (45 min).} The RHLP model, which is a mixture of experts with linear (polynomial) experts and a softmax gate. Read it with one cluster, i.e.\ ignore the outer mixture. Things to notice:
  \begin{itemize}
    \item the E-step (eq.~54) needs no forward--backward pass;
    \item the experts are updated by weighted least squares;
    \item the gate is updated by a weighted multinomial logistic regression, solved by IRLS (iteratively reweighted least squares, which is Newton's method for logistic regression).
  \end{itemize}
  This section is almost exactly Exercise A1.
\item \textbf{Bishop, as a companion (45 min).}
  \begin{itemize}
    \item \S14.5.1: mixtures of linear regression models and their EM.
    \item \S\S4.3.3--4.3.4: IRLS and multiclass logistic regression, needed for A1(g).
  \end{itemize}
\item \textbf{For Exercise A2.}
  \begin{itemize}
    \item \S4.2.1--4.2.2, pp.\ 41--45: the mixture of HMM regressions. The E-step uses forward--backward, and eqs.~39--47 give weighted updates for the initial distribution, the transition matrix, $\beta$ and $\sigma^2$.
    \item Bishop \S13.2.1--13.2.2 and \S13.2.4: the forward--backward algorithm and its scaling factors.
  \end{itemize}
\end{enumerate}
\textbf{Skip:} \S3.3--3.6 (regularised mixtures and experiments), \S4.1 (piecewise regression) and \S5 (discriminant analysis).

\begin{center}
\renewcommand{\arraystretch}{1.2}
\begin{tabular}{@{}ll@{}}
\toprule
Chamroukhi \& Nguyen & These notes \\
\midrule
$\alpha_k$ (constant mixing proportions) & $\pi_k$ with no gate input \\
$\pi_{kr}(x_j;w_k)$ (logistic gate) & $\pi_k(x)$ \\
$\pi_{kr}$ in \S4.2 & HMM initial distribution $\rho_k$ \\
$A_{k\ell r}$ & transition matrix $A_{jk}$ \\
$\tau_{ik}$, $\gamma_{ijr}$, $\xi_{ij\ell r}$ & responsibilities, regime posteriors, pairwise posteriors \\
$\beta_k$, $\sigma_k^2$ & expert weights $\beta_k$ and expert noise $s_k^2$ \\
\bottomrule
\end{tabular}
\end{center}

\subsection{Exercise A1: softmax-gated MoE with linear-regression experts}

\begin{exercise}{Exercise A1}
\textbf{Setup.} The data are $\mathcal D=\{(x_n,y_n)\}_{n=1}^N$ with $x_n\in\R^d$ (first entry $1$, the intercept) and $y_n\in\R$. There are $K$ experts and a one-hot latent $z_n$.
\begin{align*}
p(z_{nk}=1\mid x_n;w)&=\pi_k(x_n;w)=\frac{\exp(w_k\T x_n)}{\sum_{j=1}^K\exp(w_j\T x_n)},\qquad w_K=0,\\
p(y_n\mid x_n,z_{nk}=1;\beta_k,s_k^2)&=\N(y_n;\beta_k\T x_n,s_k^2),\\
\Psi&=(w_1,\dots,w_{K-1},\ \beta_1,\dots,\beta_K,\ s_1^2,\dots,s_K^2).
\end{align*}
\begin{enumerate}[label=(\alph*)]
\item Write the generative story (ancestral sampling for one data point). Draw the graphical model and mark what is observed, what is latent, and where the parameters enter.
\item Write the complete-data likelihood $p(Y,Z\mid X;\Psi)$ using 1-of-$K$ exponents, then its logarithm.
\item Write the observed-data log-likelihood $\log p(Y\mid X;\Psi)$, showing the step where $Z$ is summed out. Why is there no closed-form maximiser?
\item E-step: derive $\tau_{nk}^{(q)}=p(z_{nk}=1\mid y_n,x_n;\Psi^{(q)})$ using Bayes' rule.
\item Write $Q(\Psi;\Psi^{(q)})$. Show that it splits into one gate term that depends only on $w$, plus $K$ expert terms, each depending only on $(\beta_k,s_k^2)$.
\item Expert M-step: derive $\beta_k^{(q+1)}$ and $s_k^{2\,(q+1)}$ in closed form. Write $\beta_k$ in matrix form using the design matrix $X$, the target vector $y$ and a diagonal weight matrix.
\item Gate M-step: write the gate objective, its gradient with respect to $w_k$, and the blocks of its Hessian. Show that the objective is concave, explain why there is no closed form, and write one Newton--Raphson (IRLS) step.
\item Special case: the gate has no input, so $\pi_k$ is a constant $\alpha_k$. Show that the M-step becomes $\alpha_k^{(q+1)}=\frac1N\sum_n\tau_{nk}^{(q)}$.
\item \textbf{The panel version.} Days $t$, stocks $i\in\mathcal S_t$, one latent $z_t$ per date shared by every stock on that date, a gate that depends only on a date-level vector $u_t$, and experts
\[
p(r_{i,t+1}\mid x_{i,t},z_{tk}=1)=\N(r_{i,t+1};\beta_k\T x_{i,t},s_k^2).
\]
Redo (b), (c), (d) and (f). What changes in $\tau$, and why does it become much sharper as $N_t$ grows? (Hint: Chamroukhi \& Nguyen \S3.)
\item \emph{(Optional, but it is the key to Section 4.)} Show that
\begin{gather*}
\log L(\Psi)=Q(\Psi;\Psi^{(q)})+H(\Psi;\Psi^{(q)}),\\
H(\Psi;\Psi^{(q)})=-\E_{Z\sim p(\cdot\mid Y,X;\Psi^{(q)})}\big[\log p(Z\mid Y,X;\Psi)\big],
\end{gather*}
and use this to argue that EM never decreases the log-likelihood.
\end{enumerate}
\textbf{Solution:} Appendix A.
\end{exercise}

\subsection{Exercise A2: the same linear MoE with an HMM gate}

Exercise A2 makes three changes to A1: the regime is shared by all stocks on a day, it follows a Markov chain, and it also generates the market series that the gate sees. This is the linear version of our model. Section 5 explains the step from softmax to HMM.

\begin{exercise}{Exercise A2}
\textbf{Setup.} Days $t=1,\dots,T$, stocks $i\in\mathcal S_t$, market series $m_t$, and regime $z_t\in\{1,\dots,K\}$.
\begin{align*}
p(z_1=k)&=\rho_k, \qquad p(z_{t+1}=k\mid z_t=j)=A_{jk},\\
p(m_t\mid z_t=k)&=\N(m_t;\nu_k,\sigma_k^2),\\
p(r_{i,t+1}\mid x_{i,t},z_{t+1}=k)&=\N(r_{i,t+1};\beta_k\T x_{i,t},s_k^2),\quad\text{independently over } i\in\mathcal S_t.
\end{align*}
Tomorrow's return $r_{i,t+1}$ is predicted from today's characteristics $x_{i,t}$ and generated by tomorrow's regime $z_{t+1}$, which also generates tomorrow's market return $m_{t+1}$.
\begin{enumerate}[label=(\alph*)]
\item Draw the graphical model (a chain over $z_t$; a plate over stocks). List what is observed, what is latent, what are parameters, and what is only conditioned on.
\item Starting from the product rule and the chain rule, derive the joint $p(z_{1:T},m_{1:T},R\mid X)$, where $R$ collects all $r_{i,t+1}$. State every conditional-independence assumption you use.
\item Define the day-$(t{+}1)$ emission $e_{t+1}(k)=p(m_{t+1},\{r_{i,t+1}\}_{i\in\mathcal S_t}\mid z_{t+1}=k,X)$ (for $t=0$ take $e_1(k)=p(m_1\mid z_1=k)$). Write the observed-data likelihood as a sum over regime paths. Write the forward recursion (scaled, as in Bishop \S13.2.4) and show that it computes the likelihood in $O(TK^2)$ operations.
\item Write the backward recursion and the posteriors $\gamma_t(k)=p(z_t=k\mid\text{all data})$ and $\xi_t(j,k)=p(z_t=j,z_{t+1}=k\mid\text{all data})$.
\item Write $Q$ and derive the M-step for $\rho$, $A$, $\nu_k$, $\sigma_k^2$, $\beta_k$ and $s_k^2$. Compare $\beta_k$ with Chamroukhi \& Nguyen eq.~(46).
\item \textbf{The two-stage version.} First fit $(\rho,A,\nu,\sigma)$ by Baum--Welch on $m_{1:T}$ alone and freeze them. Then fit $(\beta,s)$ using the market-only predicted probabilities $\pi_k(t)=\xi_{t+1\mid t}(k)$ as a fixed prior for the day-$(t{+}1)$ cross-section. Write the Stage-2 E-step and M-step. Which information does the joint fit use that the two-stage fit ignores?
\item At the close of day $t$ you want to forecast $r_{i,t+1}$. Which regime probability may you use? Show that using the smoothed posterior $\gamma_{t+1}$ would be look-ahead.
\item Map the model onto Chamroukhi \& Nguyen's mixture of HMM regressions (\S4.2): what corresponds to their $n$, $K$, $R_k$, $m_i$, $y_{ij}$ and $W_{ikr}$?
\end{enumerate}
\textbf{Solution:} Appendix B.
\end{exercise}
\newpage
